% Source: https://artofproblemsolving.com/community/c3h2774733_prime_numbers_as_outputs_from_polynomial
% Stable identity: aops-topic-2774733-q01
Suppose $P$ is non-constant with integer coefficients and $P(n)$ is prime for every integer $n$. Let $p=P(0)$, a prime. For every integer $k$, $P(kp)\equiv P(0)=p\equiv0\pmod p$ (each non-constant term of $P(kp)$ is a multiple of $p$). Since $P(kp)$ is prime and divisible by $p$, it equals $p$ for every $k$. But a non-constant polynomial takes each value only finitely often (the polynomial $P(x)-p$ has finitely many roots), a contradiction. $\blacksquare$

(Constant polynomials such as $P(x)=7$ do take only prime values, which is why ``non-constant'' is needed.)
